\section{Introduction}\label{sec:intro}

Large language model (LLM) coding agents work by reading. They open logs, configuration files, stack traces and tickets through tools, and whatever a tool returns enters the model's context, where it may be stored by the provider or steered elsewhere by an injected prompt \cite{greshake2023injection,debenedetti2024agentdojo}. Workspace text routinely carries personal data and credentials \cite{meli2019git}. Agents are therefore given redaction tools, increasingly through the Model Context Protocol (MCP) \cite{mcpspec}, to check or redact a file before reading it.

Such a tool sits on the trust boundary it protects: its replies enter the model's context like any tool output, so the question is not only how much it detects but \emph{what an agent learns through it}. We identify three channels. First, the redacted text still contains every value the detectors miss. Second, the tool's ``content-free'' replies (findings with types and offsets, explanations of a decision, reports in standard formats) are themselves text, and any field that echoes input can carry a value across. Third, cost: a guardrail that adds noticeable latency to every tool call will be switched off or bypassed.

We study these channels in an open-source, on-device guardrail, SafeStream-Redactor, which exposes the same detection engine as an MCP server, a REST API and SARIF reports. We treat its pattern tiers as a fixed baseline and evaluate three other parts: a regex-gated neural name tier, the agent-facing tool surface and its per-call cost. The study was pre-registered: hypotheses, systems, corpora, metrics and decision rules were committed before any system was scored, and every departure is reported (Section~\ref{sec:method}).

\paragraph{Contributions}
\begin{enumerate}
\item \textbf{A differential egress audit for content-free tool replies.} We plant values, call every tool and search every reply in five encodings, counting a hit only if it disappears when the values are masked in the input. On the released code, the explanation tool returned case-folded person names through detector metadata, in 701 of 2,000 generated tool outputs and 24 of 129 real WNUT-17 documents; an allow-list of value-free keys removed every hit (Section~\ref{sec:egress}).
\item \textbf{A measured account of a gated neural tier.} On real user-generated text with human gold (WNUT-17), the tier raised person-name character recall over the pattern tiers by $+0.482$ [$+0.419$, $+0.548$]. The regex gate that keeps it cheap cost $0.225$ of recall against running the model on every line; batched INT8 inference broke streaming equivalence; and merged windows of up to 613 characters exceeded the declared 256-character context (Section~\ref{sec:nn}).
\item \textbf{The same audit applied to four third-party MCP redaction tools.} Three expose a restore call the agent can make, returning every hidden value; one returns all original values in the anonymiser's reply; two contradict their documentation (Section~\ref{sec:thirdparty}). A real agent, given a ticket with a prompt injection, called those restore tools and took the values back, and a pre-registered replication tested this on a second model and wording (Section~\ref{sec:agent}).
\item \textbf{The per-call cost of the guardrail by transport and redaction mode}, including a 40\,ms stall on every kept-alive REST request caused by Nagle's algorithm, found by this measurement and fixed (Section~\ref{sec:cost}).
\end{enumerate}
Code, corpora loaders, frozen manifests, raw results and the protocol are public (\repolink).
